\begin{problem}{Small Multipliers}
Indian computational manuals reduce a planet's revolution count and the civil-day count of a great cycle to a coprime pair, then expand the ratio as a continued fraction to obtain multipliers and divisors small enough to compute with by hand. In a great cycle of \(1{,}577{,}917{,}500\) civil days the Moon completes \(57{,}753{,}336\) revolutions and its apogee completes \(488{,}219\); the difference is the number of revolutions of the lunar anomaly.

Form the ratio of civil days to revolutions of the anomaly, reduce it to lowest terms, and expand it as a simple continued fraction. Taking the convergents in order and counting each convergent once, find the sum of all convergent denominators below \(10^6\). Note that the first two convergents both have denominator \(1\) and both count.
\end{problem}